
\subsubsection{决策变量与约束}

决策变量 $x_r \in \{0,1\}$，$y^{O}_{r,f}$、$y^{I}_{r,f}$、$y^{S}_{r,(f_1,f_2)} \in \mathbb{Z}_{\geq 0}$。约束条件包括：

\textbf{需求覆盖约束}（按日）：$\sum_{r} y^{O}_{r,f} = n^{O}_{d,f}$，$\sum_{r} y^{I}_{r,f} = n^{I}_{d,f}$，$\sum_{r} y^{S}_{r,(f_1,f_2)} = n^{S}_{d,(f_1,f_2)}$。

\textbf{航段级容量约束}：同问题二公式~\eqref{eq:q2_capacity}，原样迁移。

\textbf{时间窗约束}：每条路径起飞时刻 $t^{dep}_r$ 取固定值 $\in [E_r, L_r]$，需求-路径兼容性在路径生成阶段预计算（命题~\ref{prop:cpm}）。

\textbf{机队约束}：
\begin{equation}
\sum_{\substack{r \in \mathcal{R}^{a,t}_d \\ \tau \in [t^{dep}_r,\; t^{ret}_r + 30]}} x_r \leq |\mathcal{M}_{a,t}|
\label{eq:q3_fleet}
\end{equation}

\subsubsection{两阶段求解策略}

题目要求两阶段求解：阶段一在不考虑临时任务的情况下编排（3840 人），得到总飞机使用时间 $T^*$；阶段二在 $T \leq T^*$ 约束下插入临时任务（160 人）。求解流程如图~\ref{fig:q3_flow} 所示。

% -*- coding: utf-8 -*-
% 问题三两阶段求解流程图（服务日入列 + 事件驱动时刻 + 覆盖兜底 + 逐日阶段二）
\begin{figure}[htbp]
\centering
\begin{tikzpicture}[
    box/.style={draw, rounded corners=3pt, minimum width=3.4cm, minimum height=0.75cm,
        align=center, font=\footnotesize, fill=blue!8},
    boxg/.style={draw, rounded corners=3pt, minimum width=3.4cm, minimum height=0.75cm,
        align=center, font=\footnotesize, fill=green!10},
    boxo/.style={draw, rounded corners=3pt, minimum width=3.4cm, minimum height=0.75cm,
        align=center, font=\footnotesize, fill=orange!12},
    boxr/.style={draw, rounded corners=3pt, minimum width=3.4cm, minimum height=0.75cm,
        align=center, font=\footnotesize, fill=red!8},
    boxp/.style={draw, rounded corners=3pt, minimum width=3.4cm, minimum height=0.75cm,
        align=center, font=\footnotesize, fill=purple!10},
    arrow/.style={->, >=stealth, thick, draw=gray!70},
    darrow/.style={->, >=stealth, thick, draw=gray!70, dashed},
]

% 阶段一标题
\node[font=\small\bfseries, fill=green!15, draw, rounded corners, minimum width=12.5cm, minimum height=0.5cm] (s1) at (0, 4.6) {阶段一：固定排班（3840人）};

% 第一行：预处理与列生成（4 列）
\node[boxg] (window) at (-4.6, 3.5) {时间窗解析\\可行日期集 $D_p$（服务日自由）};
\node[boxg] (cols) at (0, 3.5) {全周候选列池\\事件驱动时刻 $\times$ (机场,机型) $\times$ 多设施序列};
\node[boxg] (load) at (4.6, 3.5) {EDF 双序满员装载\\mandatory 优先 + 临时填残座};
\node[boxp] (cover) at (8.2, 3.5) {覆盖兜底\\专属直达列族};

% 第二行：剪枝与 MILP（3 列）
\node[box] (prune) at (-3.4, 2.2) {列池剪枝 \\(40,10) $\to$ 约 7 万列};
\node[box] (mip) at (1.2, 2.2) {阶段一 COPT MILP\\$\min T$，3840 全覆盖 + 机队并发};
\node[boxg] (assign1) at (5.8, 2.2) {飞机指派\\区间图并发保证};

% 阶段一输出
\node[boxg] (tstar) at (0, 1.0) {每日 $T^*_d$（日预算）$\Rightarrow T^*$};

% 阶段二标题
\node[font=\small\bfseries, fill=orange!15, draw, rounded corners, minimum width=12.5cm, minimum height=0.5cm] (s2) at (0, -0.6) {阶段二：临时任务加入（160人，$T \le T^*$）};

% 阶段二流程（3 列）
\node[boxo, below=0.35cm of s2, xshift=-3.2cm] (day_cols) {逐日列裁剪\\当日可行临时需求};
\node[boxo, below=0.35cm of s2] (prize) {逐日 prize-collecting MIP\\$\max \sum z_p$，日预算 $T^*_d$};
\node[boxo, below=0.35cm of s2, xshift=3.2cm] (rerun) {阶段一架次作保底列\\MIP start 防无解};

\node[boxr, below=0.55cm of prize] (output) {输出 q3-routes / q3-assignments\\（3840 全覆盖，临时按预算承运）};

% 主箭头
\draw[arrow] (window) -- (cols);
\draw[arrow] (cols) -- (load);
\draw[arrow] (load) -- (cover);
\draw[arrow] (cover.west) -- (prune.east);
\draw[arrow] (prune) -- (mip);
\draw[arrow] (mip) -- (assign1);
\draw[arrow] (assign1.south) |- (tstar.east);
\draw[arrow] (tstar) -- (s2);
\draw[arrow] (s2) -- (day_cols);
\draw[arrow] (s2) -- (prize);
\draw[arrow] (s2) -- (rerun);
\draw[arrow] (day_cols) -- (prize);
\draw[arrow] (rerun) -- (prize);
\draw[arrow] (prize) -- (output);

% 覆盖兜底反馈
\draw[darrow] (cover.south) -- ++(0,-0.15) -| (window.south);
\node[font=\scriptsize, align=center, anchor=north] at (8.2, 2.75) {缺失人员\\补列};

\end{tikzpicture}
\caption{问题三两阶段求解流程（服务日入列 + 事件驱动时刻 + 覆盖兜底 + 逐日阶段二）}
\label{fig:q3_flow}
\end{figure}


阶段一按天解耦为 7 个独立单日问题，通过需求-日期分配（跨天需求选负载最低的可行日期，贪心均衡）后逐日求解（CPM 收缩 $\to$ 受限路径池 $\to$ COPT MIP）。每个单日问题采用两层求解：第一层单目标严格最小化当日总飞机使用时间（时限 180 s，得到当日 $T^*_d$）；第二层固定第一层的路径选择，在 $T_d = T^*_d$（无上界差）下优化该日人员在途时间与燃油的加权目标。全部 7 日汇总得阶段一总时间 $T^* = \sum_d T^*_d$。

阶段二采用临时任务背包化策略。阶段一固定排班确定后，剩余资源为``时间缝隙 + 空余座位''。将问题重构为最大覆盖/背包问题：
\begin{equation}
\max \sum_{p \in \mathcal{P}^{*}} z_p \quad \text{s.t.} \quad \sum_{p} \ell_{r,p} \cdot z_p \leq \Delta_r, \quad \sum_r t_r \leq T^*
\label{eq:q3_knapsack}
\end{equation}

优先插入已有架次空余座位（零新增时间），不足时在 $T^*$ 预算内新增架次。此策略避免重跑大 MIP，且可给出上界；临时任务的插入始终受 $T \le T^*$ 硬约束，不破坏阶段一的总时间最优性。

\subsubsection{硬约束兜底机制}

为保证最终 \texttt{q*-assignments.csv} 中不存在任何违规分配（航段超载、时间窗越界、日期错配），在求解器输出与人员写入之间设置三层硬约束兜底：

\textbf{(1) 容量保护。}问题一中以 ``路线总已用人数 + 设施配额'' 双重限流（\texttt{route\_used + per\_fac\_limit}），任一架次在机人数不超过座位数；问题二将容量细化到航段级，维护每架次各航段剩余座位池 \texttt{seg\_remaining}，出海/海返/穿梭人员在共享池中按区间扣减，保证任意航段不超载；问题三在写入阶段对每条分配调用 \texttt{seg\_has\_space} 逐航段校验（穿梭还需满足上机位置先于下机位置的方向约束），顺序分配导致容量变化时还会二次校验。

\textbf{(2) 时间窗保护。}问题三对每条需求计算 \texttt{check\_compatibility}：航班在登机停靠点的起飞时刻不早于需求最早可接人时刻、在下机停靠点的到达时刻不晚于需求最晚到达时刻，二者同时满足才允许分配。

\textbf{(3) 日期保护。}需求-日期分配仅给出建议日期；写入阶段对未匹配需求执行二次匹配，逐条校验航班日期属于 \texttt{get\_feasible\_days} 返回的可行日期集合，禁止跨可行日期的搭错航班。当某日 MIP 无可行解时，\texttt{failed\_day\_demands} 机制将该日需求按可行日期与负载均衡重分配到已求解日，仍无法满足的如实标记为未服务。

该兜底机制保证验证脚本（V1--V10）中容量、时间窗、日期、OD 一致性、方向顺序、重复分配与航油七类硬约束全部零违规，是结果可信性的最后防线。

\subsection{求解结果}

问题三的求解结果如表~\ref{tab:q3_result} 所示。

\begin{longtable}{lc}
\caption{问题三求解结果} \label{tab:q3_result} \\
\toprule
指标 & 数值 \\
\midrule
\endfirsthead
\toprule
指标 & 数值 \\
\midrule
\endhead
\bottomrule
\endfoot
\bottomrule
\endlastfoot
总飞机使用时间 & 35451 min（590.85 h） \\
人员总在途时间 & 486738 min（8112.30 h） \\
总架次数 & 207 \\
总燃油消耗量 & 293456.7 kg \\
座位利用率 & 35.81\% \\
临时任务满足率 & 99.38\% \\
\end{longtable}
